%***************************************************************************************************
% preamble/macros.tex - makra semantyczne, skróty matematyczne, twierdzenia, listy.
%***************************************************************************************************

% ---------------------------------------------------------------------------------------------
% Makra semantyczne (spójny wygląd kodu, plików, poleceń i skrótów)
% \detokenize pozwala wpisywać znaki specjalne (np. _ oraz %) bez zabezpieczania.
% ---------------------------------------------------------------------------------------------
\newcommand{\code}[1]{\texttt{\detokenize{#1}}}   % fragment kodu / wyrażenie
\newcommand{\file}[1]{\texttt{\detokenize{#1}}}   % nazwa pliku lub ścieżka
\newcommand{\cmd}[1]{\texttt{\detokenize{#1}}}    % polecenie konsoli
% Polecenie LaTeX-a w tekście: \latexcmd{section} wypisuje \section; nawiasy klamrowe zapisuj jako \{ \}
\newcommand{\latexcmd}[1]{\texttt{\textbackslash #1}}
% Skrót kryptograficzny lub inny długi ciąg znaków - zawija się w dowolnym miejscu
\newcommand{\hash}[1]{{\ttfamily\footnotesize\seqsplit{#1}}}

% Pogrubiony nagłówek wewnątrz sekcji (zamiast ręcznego \medskip\noindent\textbf{..}\\)
\newcommand{\lead}[1]{\par\medskip\noindent\textbf{#1}\par\nopagebreak\smallskip}

% Ramka z podpowiedzią w przykładowych rozdziałach - usuń wywołania \TemplateNote przy pisaniu własnej treści
\newcommand{\TemplateNote}[1]{%
  \par\medskip\noindent
  \fcolorbox{lgray}{gray!8}{%
    \parbox{\dimexpr\linewidth-2\fboxsep-2\fboxrule\relax}{\small\textbf{Wskazówka szablonu.} #1}}%
  \par\medskip
}

% Rozdział bez numeru (dopisany do spisu treści): \UnnumberedChapter{Streszczenie}
\newcommand{\UnnumberedChapter}[1]{%
  \clearpage\phantomsection
  \addcontentsline{toc}{chapter}{#1}%
  \chapter*{#1}%
}

% Pytanie i odpowiedź: \begin{qa}{Treść pytania} Odpowiedź... \end{qa}
\newcounter{qacount}
\newenvironment{qa}[1]{%
  \refstepcounter{qacount}%
  \par\medskip\noindent\textbf{Pytanie nr~\theqacount}\par\nopagebreak
  \noindent\textit{\enquote{#1}}\par\nopagebreak\smallskip
  \noindent\textbf{Odpowiedź}\par\nopagebreak
}{\par\medskip}

% Spis skrótów: \begin{abbreviations} \item[HTTP] Hypertext Transfer Protocol ... \end{abbreviations}
\newenvironment{abbreviations}
  {\begin{description}[style=multiline, leftmargin=3.5cm, font=\bfseries, itemsep=2pt]}
  {\end{description}}

% Tylda wyśrodkowana w pionie (przybliżenie)
\newcommand{\textapprox}{\raisebox{0.5ex}{\texttildelow}}

% Znaczniki w tabelach (\cmark, \xmark) i pogrubiona komórka nagłówka
\newcommand{\cmark}{\ding{51}}
\newcommand{\xmark}{\ding{55}}
\newcommand{\hcell}[1]{\textbf{#1}}

% ---------------------------------------------------------------------------------------------
% Skróty matematyczne
% ---------------------------------------------------------------------------------------------
\newcommand{\RR}{\mathbb{R}}
\newcommand{\NN}{\mathbb{N}}
\newcommand{\QQ}{\mathbb{Q}}
\newcommand{\ZZ}{\mathbb{Z}}
\newcommand{\IMP}{\rightarrow}
\newcommand{\IFF}{\leftrightarrow}

% ---------------------------------------------------------------------------------------------
% Twierdzenia, definicje, przykłady, uwagi (numeracja w obrębie rozdziału, wspólny licznik)
% ---------------------------------------------------------------------------------------------
\theoremstyle{plain}
\newtheorem{theorem}{Twierdzenie}[chapter]
\newtheorem{lemma}[theorem]{Lemat}
\newtheorem{corollary}[theorem]{Wniosek}
\theoremstyle{definition}
\newtheorem{definition}[theorem]{Definicja}
\newtheorem{example}[theorem]{Przykład}
\theoremstyle{remark}
\newtheorem*{remark}{Uwaga}
\renewcommand{\qedsymbol}{$\blacklozenge$}

% ---------------------------------------------------------------------------------------------
% Listy (zwykłe itemize / enumerate - jeden spójny styl w całym dokumencie)
% ---------------------------------------------------------------------------------------------
\setlist{topsep=0.4em, itemsep=0.2em, parsep=0pt, partopsep=0pt}
\setlist[itemize,1]{label=\textbullet, leftmargin=*}
\setlist[itemize,2]{label=\textendash, leftmargin=*}
\setlist[itemize,3]{label=$\diamond$, leftmargin=*}
\setlist[enumerate,1]{label=\arabic*., leftmargin=*}
\setlist[enumerate,2]{label=\alph*), leftmargin=*}

% Mniej "rozjeżdżonych" akapitów w wąskich kolumnach i listach
\setlength{\emergencystretch}{2em}
